\lecture{7.1}{Limitations of Naive G-Computation and IPTW}
\label{lec:7-1}

\subsection{Dependence on Nuisance Regression Consistency}

The point-treatment ATE is identified, under consistency, conditional exchangeability,
and positivity, by
\[
\ATE=\E\{\Qbar(1,W)-\Qbar(0,W)\},
\qquad
\Qbar(a,w)=\E(Y\given A=a,W=w),
\]
and also by inverse probability weighting:
\[
\E\{\Ya{a}\}
=
\E\left\{\frac{\ind(A=a)}{g_a(W)}Y\right\},
\qquad
g_a(w)=\Prob(A=a\given W=w).
\]
These formulas lead directly to the estimators studied in Chapter 4:
\[
\psi_{n,G}(a)=\frac{1}{n}\sum_{i=1}^n\Qbar_n(a,W_i)
\]
and
\[
\psi_{n,\IPTW}(a)=
\frac{1}{n}\sum_{i=1}^n
\frac{\ind(A_i=a)}{g_{n,a}(W_i)}Y_i.
\]

Each estimator relies on consistent estimation of a nuisance regression. The
g-computation estimator relies on $\Qbar_n$ being close to $\Qbar$. The IPTW
estimator relies on $g_{n,a}$ being close to $g_a$. If the required nuisance function is
estimated poorly, the resulting causal effect estimator can be biased even in large
samples.

\begin{assumptionbox}
All estimators in this chapter require the same causal identification assumptions as
before:
\[
Y=Y(A),\qquad \Ya{a}\ci A\given W,\qquad 0<\Prob(A=a\given W=w)<1.
\]
Doubly robust estimation improves statistical robustness to nuisance-model
misspecification; it does not remove the need for these causal assumptions.
\end{assumptionbox}

\subsection{Outcome Regression Misspecification}

Suppose the analyst fits a main-terms outcome regression
\[
\Qbar_n(a,w)\approx \beta_0+\beta_1a+\beta_2w,
\]
but the true conditional mean contains an interaction:
\[
\Qbar(a,w)=\beta_0+\beta_1a+\beta_2w+\beta_3aw.
\]
Then $\Qbar_n(1,w)-\Qbar_n(0,w)$ may be systematically wrong for some values of
$w$. Since g-computation averages these predicted contrasts over the covariate
distribution, the resulting ATE estimator can converge to the wrong target.

This is not a small-sample problem. With enough data, the fitted main-terms model may
converge very precisely to the best linear approximation within the misspecified
model class. Precision around the wrong function is still bias for the causal
estimand.

\subsection{Propensity Score Misspecification}

IPTW has a parallel vulnerability. If the propensity score model is misspecified, the
weights
\[
\frac{\ind(A_i=a)}{g_{n,a}(W_i)}
\]
do not create the intended pseudo-population. Some covariate strata may remain
overrepresented and others underrepresented after weighting. In extreme cases, an
incorrect propensity model can also create highly variable weights without correcting
confounding.

For example, if the true propensity score is nonlinear,
\[
g_1(w)=\expit(\alpha_0+\alpha_1w+\alpha_2w^2),
\]
but the analyst fits $\expit(\gamma_0+\gamma_1w)$, the fitted probabilities can be
wrong exactly in regions where treatment assignment is most imbalanced.

\subsection{Flexible Regression Complications}

Super Learning and other flexible regression methods help reduce the risk of severe
misspecification. They allow $\Qbar_n$ and $g_n$ to adapt to nonlinearities,
interactions, and high-dimensional predictor structures.

However, flexible regression creates two new issues. First, flexible learners may
converge more slowly than parametric estimators. Second, their sampling distributions
can be difficult to characterize. The final causal estimator is no longer a simple
average of independent terms with fixed nuisance functions; it contains data-adaptive
estimated functions.

\begin{figure}[H]
\centering
\begin{tikzpicture}[
  node distance=1.15cm,
  box/.style={draw, rounded corners=2pt, minimum width=3.1cm, minimum height=.75cm},
  arr/.style={-{Latex[length=2.5mm]}, thick}
]
\node[box] (data) {Observed data};
\node[box, below left=1cm and .9cm of data] (q) {Estimate\\ $\Qbar_n$};
\node[box, below right=1cm and .9cm of data] (g) {Estimate\\ $g_n$};
\node[box, below=1cm of q] (gc) {G-computation\\ depends on $\Qbar_n$};
\node[box, below=1cm of g] (ipw) {IPTW\\ depends on $g_n$};
\node[box, below right=1cm and .1cm of gc] (aipw) {Doubly robust\\ estimators use both};
\draw[arr] (data) -- (q);
\draw[arr] (data) -- (g);
\draw[arr] (q) -- (gc);
\draw[arr] (g) -- (ipw);
\draw[arr] (q) -- (aipw);
\draw[arr] (g) -- (aipw);
\end{tikzpicture}
\caption{Naive g-computation and IPTW each depend on one nuisance function; doubly robust estimators combine both.}
\label{fig:ch7-nuisance-dependence}
\end{figure}

\subsection{Lack of Standard Inference Basis}

For a simple sample mean, the central limit theorem gives a standard basis for
confidence intervals. For a parametric regression coefficient, classical regression
theory gives approximate standard errors under model assumptions. For naive
g-computation or IPTW with highly adaptive nuisance estimates, neither argument
directly applies without additional work.

The statistical question is not only whether $\psi_n$ is consistent. We also need an
approximation of the form
\[
\sqrt n(\psi_n-\psi)\overset{d}{\longrightarrow}N(0,\sigma^2),
\]
and a consistent estimator of $\sigma^2$. Doubly robust estimators are designed to
have an asymptotic linear representation under appropriate conditions, which gives a
direct route to inference.

\subsection{Bootstrap Limitations}

The nonparametric bootstrap is often useful, but it is not a universal remedy. With
highly adaptive learners, variable selection, non-smooth tuning, or near-positivity
violations, bootstrap distributions can fail to approximate the true sampling
distribution. The bootstrap may also be computationally expensive because every
bootstrap replicate must refit all nuisance functions.

\begin{keyidea}
The limitation of naive g-computation and IPTW is not that they are wrong as
identification formulas. The limitation is that, as estimators, they rely heavily on one
nuisance regression and do not automatically provide stable, valid inference when
that nuisance regression is learned flexibly.
\end{keyidea}
